\section{Conclusion}
\label{sec:conclusion}

This paper presented a systematic literature review of 54 studies on collaborative SLM--LLM systems across hybrid, multi-agent, ensemble, and retrieval-centered architectures, with monolithic LLMs treated as baselines rather than as a primary study class. The four research questions yield consistent, actionable answers.

\textbf{RQ1 (Performance).} Hybrid and multi-agent architectures consistently outperform monolithic SLM baselines and can match or exceed LLM baselines in narrower, structured settings. Performance gains are driven by orchestration strategy---selective escalation, role decomposition, verification loops---rather than by model size alone, and are strongest in domain-specific tasks with well-defined subtask boundaries.

\textbf{RQ2 (Efficiency).} The strongest efficiency gains, often exceeding 80\% cost reduction, arise from token-level routing and selective escalation. Latency effects are polarized: local SLM inference is fast, but orchestration layers frequently introduce overhead that negates throughput gains. Energy remains the least reported metric; available evidence shows that generation length and reasoning depth matter as much as model size for per-query energy use, and pushing small models into long Chain-of-Thought traces can exceed the per-query energy cost of a larger model prompted more directly~\cite{energytoken2025}.

\textbf{RQ3 (Orchestration).} Orchestration mechanisms cluster into five recurring families: token-level verification, query-level routing, RAG, task decomposition, and multimodal fusion. The choice among them is determined primarily by deployment constraints and data modalities rather than by raw model capacity.

\textbf{RQ4 (Domain Suitability).} Domain-specific constraints---privacy mandates in healthcare, latency budgets in telecom, hardware limits at the edge, and syntax rigidity in code tasks---are the primary drivers of architecture selection. Collaborative SLM--LLM systems deliver their strongest results when the workflow can be partitioned to align each component with the constraint it best addresses.

Across all four questions, a consistent principle emerges: the presence of an explicit control layer---router, planner, retriever, or fusion module---is the most reliable differentiator between high- and low-performing collaborative designs (39 of 54 included studies). Strong results are reliably associated with these control layers rather than with model scale alone.

The literature is stronger in architectural innovation than in evaluation discipline. Many studies report accuracy improvements, but far fewer report latency, cost, or energy in a form that supports fair comparison. The current evidence is sufficient to show that hybridization \textit{can} work, but often insufficient to show \textit{when} it is the better deployment choice under realistic resource constraints. The strongest positive results remain concentrated in narrower or more structured settings such as telecom question answering, selected healthcare tasks, optimisation, and industrial monitoring. Energy evidence is especially limited: direct measurement is rare, and proxy improvements such as reduced computational redundancy~\cite{resiliomate2025} do not establish verified energy savings. Future work should treat energy as a first-class evaluation outcome and report energy-per-token, reasoning length, hardware configuration, and escalation behavior under comparable conditions. Standardized benchmarks for collaborative SLM--LLM systems would further reduce the current reliance on paper-specific evaluation protocols.

The full comparative analysis table and data extraction spreadsheet supporting this review are available at \url{https://bit.ly/SLM-LLM-Data-2026}.
